\color{unimain}{\subsection{Газовая защита РПН}}
\color{black}

\begin{enumerate}

\item
Газовая защита устройства РПН предназначена для защиты от повреждений внутри бака РПН трансформатора, сопровождающихся выделением газа.   

\item
Алгоритм работы логики ГЗ РПН приведен на рисунке \ref{pict:GZ_RPN}, параметры для настройки ГЗ РПН приведены в таблице \ref{app_set:gzrpn}.

\begin{figure}[H]
\centering
\includegraphics[width=0.8\textwidth,height=0.8\textheight,keepaspectratio]{\fbpath/07.01 РЗТ/011.1073 ГЗ РПН/img/GZ_RPN.pdf}
\captionof{figure}{Функциональная схема логики ГЗ РПН}\label{pict:GZ_RPN}
\end{figure}

\item
Предусмотрена возможность приема сигналов срабатывания ГЗ РПН и низкой изоляции цепи ГЗ РПН, по разным цепям от двух источников ПДС1 и ПДС2. Для приема сигналов от ПДС1 предусмотрены входные сигналы \textit{«Ср. ГЗ РПН 1 цепь»} и \textit{«Низкая изол. ГЗ РПН 1 цепь»}, а для приема сигналов от ПДС2 предусмотрены входные сигналы \textit{«Ср. ГЗ РПН 2 цепь»} и \textit{«Низкая изол. ГЗ РПН 2 цепь»}. Если сигналы срабатывания ГЗ РПН заводятся непосредственно со шкафа ШОМ силового трансформатора конфигурируются только входные сигналы \textit{«Ср. ГЗ РПН 1 цепь»} и \textit{«Низкая изол. ГЗ РПН 1 цепь»}, а действие ГЗ РПН по второй цепи выводиться из работы с помощью программного переключателя \textbf{«ГЗ РПН 2 цепь»}.

\item
Ввод в работу ГЗ РПН осуществляется с помощью программных переключателей \textbf{«ГЗ РПН 1 цепь»} и \textbf{«ГЗ РПН 2 цепь»}. Оперативный перевод ГЗ РПН на сигнал производится с помощью команды управления \textit{«Перевод ГЗ РПН на сигнал»}.

\item
Оперативный ввод/вывод ГЗ РПН производится с помощью БУ (блока управления функциями) \textit{«Вывод ГЗ РПН»}.

\item
Непрерывный контроль изоляции цепи ГЗ РПН осуществляется с помощью реле контроля изоляции (РКИ). В случае снижения сопротивления изоляции РКИ подает сигнал \textit{«Низкая изол. ГЗ РПН 1 цепь»} (\textit{«Низкая изол. ГЗ РПН 2 цепь»}) для блокировки срабатывания ГЗ РПН. Время срабатывания блокировки от сигналов РКИ задается с помощью уставки \textbf{«Тбл ГЗ РПН»}. Ввод блокировки ГЗ РПН от РКИ осуществляется с помощью программного переключателя \textbf{«Блок ГЗ РПН»}. Для сброса блокировки от РКИ, используется команда управления \textit{«Возврат ГЗ РПН»}.

\end{enumerate} 